\chapter{Developer Experience dan Portal Pengembang}

\section*{Tujuan Pembelajaran}
\begin{itemize}
  \item Menjelaskan unsur \textit{developer experience} yang menentukan cepat atau lambatnya mitra terhubung.
  \item Merancang alur \textit{onboarding} mitra dari pendaftaran sampai produksi.
  \item Menyusun kerangka portal pengembang beserta ukuran keberhasilannya.
\end{itemize}

\section{Pendahuluan}

% TODO
% Cakupan: klinik mitra butuh berbulan-bulan untuk terhubung karena dokumentasi tersebar, kredensial
% dikirim lewat surel, dan lingkungan uji tidak tersedia.

\section{Mitra sebagai Pengguna API}

% TODO
% Cakupan: persona pengguna API (vendor SIMRS klinik, tim TI asuransi, pengembang aplikasi kesehatan
% perusahaan); kebutuhan dan keterbatasan tiap persona; pengambil keputusan bisnis di pihak mitra.
% Kasus Medikara: persona klinik mitra dengan tim TI kecil.

\section{Developer Experience dan Waktu sampai Panggilan Pertama}

% TODO
% Cakupan: arti developer experience; waktu sampai panggilan pertama (time to first call) sebagai
% ukuran utama; temuan Postman 2025 tentang dokumentasi yang tidak konsisten dan API yang sulit
% ditemukan; hambatan umum mitra.

\section{Dokumentasi, Sandbox, dan SDK}

% TODO
% Cakupan: dokumentasi referensi, panduan, dan contoh; sandbox dengan data uji sintetis agar data
% pasien tidak dipakai; SDK dan koleksi contoh permintaan; menjaga dokumentasi tetap sesuai kontrak.
% Pendalaman teknis: membangkitkan halaman dokumentasi dari kontrak OpenAPI Bab 9 dengan alat
% gratis, lalu mengukur waktu sampai panggilan pertama bersama rekan yang belum mengenal API itu.

\section{Portal Pengembang dan Katalog API}

% TODO
% Cakupan: fungsi portal (katalog, dokumentasi, kredensial, status layanan, dukungan); struktur
% halaman minimum; contoh portal SNAP Developer Site dan dokumentasi SATUSEHAT sebagai pembanding.
% Kasus Medikara: kerangka portal pengembang Medikara.

\section{Alur Onboarding dan Verifikasi Mitra}

% TODO
% Cakupan: langkah pendaftaran, verifikasi identitas dan legalitas mitra, akses sandbox, uji
% kelayakan, dan akses produksi; penerbitan kredensial yang aman; titik persetujuan bisnis dan
% keamanan.
% Kasus Medikara: onboarding laboratorium mitra dan asuransi.

\section{Syarat Penggunaan, Tingkat Layanan, dan Dukungan}

% TODO
% Cakupan: syarat penggunaan dan batas pemakaian data; SLA (Service Level Agreement) dan cara
% mengukurnya; saluran dukungan dan eskalasi; komunitas pengembang; aspek kontrak dengan pihak
% ketiga dibahas rinci di IT40813.

\section{Aktivitas: Rancangan Portal Pengembang dan Onboarding Mitra}

% TODO
% Tujuan aktivitas: merancang pengalaman mitra dari pendaftaran sampai produksi.
% Peran dalam tim: anggota berlatar bisnis dan hukum menyusun syarat penggunaan, tingkat layanan,
% dan alur persetujuan; anggota berlatar teknis menyusun kerangka portal dan isi dokumentasi
% minimum.
% Langkah: pilih satu jenis mitra; susun persona; gambar alur onboarding; susun kerangka portal;
% tetapkan target waktu sampai panggilan pertama dan ukuran keberhasilan portal.
% Luaran dan tonggak capstone: rancangan portal pengembang dan alur onboarding mitra, melengkapi C3.
% Berkas sources yang direncanakan: sources/chapter11/alur-onboarding-mitra.md,
% sources/chapter11/kerangka-portal-pengembang.md, sources/chapter11/templat-syarat-penggunaan.md

\section{Ringkasan}

% TODO

\section{Latihan}

\begin{enumerate}
  \item Jelaskan mengapa dokumentasi yang lengkap belum tentu menghasilkan \textit{developer experience} yang baik.
  \item Susun lima langkah \textit{onboarding} untuk laboratorium mitra, lalu tandai langkah yang paling mungkin menjadi hambatan.
  \item Tentukan tiga ukuran keberhasilan portal pengembang dan jelaskan cara mengumpulkan datanya.
\end{enumerate}
